\section{Resumen}
El presente documento describe el proyecto de optimización de procesos en la infraestructura tecnológica del Centro de Innovación e Integración de Tecnologías Avanzadas (CIITA) de Ciudad Juárez. El CIITA, como parte del Instituto Politécnico Nacional (IPN), enfrenta una crisis operativa derivada del deterioro de su infraestructura tecnológica y un cambio en su modelo de financiamiento. Este problema se manifiesta en tres áreas críticas: (1) infraestructura tecnológica obsoleta, (2) configuración insegura de la red, y (3) gestión ineficiente del soporte técnico. Estas deficiencias no solo limitan la capacidad del centro para cumplir con sus objetivos estratégicos, sino que también aumentan los riesgos de seguridad y afectan su reputación ante inversionistas y socios estratégicos.

La justificación del proyecto se basa en la necesidad de garantizar la continuidad operativa del CIITA, cumplir con normativas como la Ley Federal de Protección de Datos Personales en Posesión de Particulares (LFPDPPP) y evitar sanciones económicas que podrían alcanzar millones de pesos. Además, se destaca la importancia de mantener certificaciones internacionales, como la ISO 27001, que exigen la implementación de medidas robustas de seguridad de la información.

El proyecto propone una serie de acciones para modernizar la infraestructura tecnológica, mejorar la configuración de la red y optimizar la gestión del soporte técnico. Estas acciones no solo buscan resolver los problemas actuales, sino también fortalecer la capacidad del CIITA para seguir siendo un motor de desarrollo tecnológico y económico en la región. La inversión requerida para estas mejoras es estratégica, considerando el impacto positivo que tendrá en la formación de capital humano, la investigación y desarrollo, y la competitividad industrial de Ciudad Juárez.

En resumen, este proyecto representa una oportunidad para asegurar la sostenibilidad del CIITA, garantizar el cumplimiento de normativas y estándares internacionales, y mantener su posición como un centro de innovación y desarrollo tecnológico de vanguardia en México.